\chapter*{Introducción}
\addcontentsline{toc}{chapter}{Introducción}
\markboth{Introducción}{Introducción}

Las ecuaciones diferenciales parciales (EDPs) elípticas aparecen de manera natural en numerosos problemas de la física, la ingeniería y las ciencias aplicadas, como la electrostática, la difusión de calor o la mecánica de fluidos. Dentro de esta familia, la ecuación de Poisson ocupa un lugar central por su amplia aplicabilidad y por ser un modelo prototípico para el estudio de métodos analíticos y numéricos. Tradicionalmente, su solución se ha abordado mediante esquemas de diferencias finitas~\cite{leveque2007}, elementos finitos~\cite{strang1973, babuska1994} o volúmenes finitos, cuya teoría está bien establecida pero que pueden enfrentar limitaciones en geometrías complejas o en altas dimensiones.

Las \emph{redes neuronales físicamente informadas} (\emph{Physics-Informed Neural Networks}, PINNs) representan la solución mediante una red y entrenan sus parámetros penalizando el residuo de la ecuación y el defecto de frontera. Esta formulación, desarrollada entre otros por Raissi et al.~\cite{raissi2017physics, raissi2019}, permite incorporar también observaciones interiores. Entre los antecedentes del uso de redes neuronales para aproximar soluciones de EDPs se encuentran Dissanayake y Phan-Thien (1994)~\cite{dissanayake1994neural} y, posteriormente, Lagaris et al. (1998)~\cite{lagaris1998artificial}; las arquitecturas y técnicas de entrenamiento actuales se apoyan en avances posteriores del aprendizaje profundo~\cite{goodfellow2016deep, lecun2015deep}.

Una pérdida residual pequeña en un conjunto finito de puntos no asegura, por sí sola, que el residuo sea pequeño en todo el dominio. Tampoco garantiza que el optimizador encuentre un minimizador con las propiedades requeridas por los resultados de convergencia. Estas dos dificultades motivan el estudio matemático y los experimentos de esta tesis.

La motivación específica de este trabajo es doble. Por un lado, se busca contar con un marco matemático que permita justificar, bajo hipótesis claras, la convergencia de las PINNs hacia la solución exacta de una EDP elíptica, tanto en normas débiles como fuertes. Por otro lado, se pretende contrastar estos resultados teóricos con experimentos numéricos concretos para la ecuación de Poisson en una y dos dimensiones, analizando el comportamiento de los errores y comparando el desempeño de las PINNs con métodos clásicos.

La tesis se organiza en cuatro capítulos de desarrollo y uno de conclusiones. En el capítulo I, dedicado a los preliminares matemáticos, se introduce el marco analítico necesario basándose en la teoría clásica de EDPs elípticas~\cite{gilbarg2001elliptic}: definición de operadores lineales, formulación del problema de Dirichlet, espacios de Hölder, espacios $L^2$ y de Sobolev, principio del máximo, celdas de Voronoi y el teorema de existencia de Schauder. Estos elementos se utilizarán posteriormente para formular y probar los resultados de convergencia.

En el capítulo II, dedicado a las redes neuronales, se presentan las nociones básicas de redes neuronales artificiales~\cite{bishop2006pattern} y se describe la construcción de las PINNs. Se detallan las funciones de activación, el algoritmo de optimización Adam~\cite{adam2015method}, el uso de diferenciación automática y la formulación de la función de pérdida, así como una versión regularizada adecuada para problemas elípticos.

El capítulo III, dedicado al análisis de convergencia, constituye el núcleo teórico de la tesis. Siguiendo la línea de trabajo de Shin et al.~\cite{shin2020convergence}, se establecen lemas, hipótesis y teoremas que describen las condiciones bajo las cuales una familia de redes neuronales minimizadoras de la pérdida PINN converge hacia la solución clásica de la ecuación de Poisson. Se analizan tanto la convergencia en la norma $C^0(U)$ como el paso a normas más fuertes, como $H^1(U)$, apoyándose en construcciones basadas en celdas de Voronoi.

Finalmente, en el capítulo IV, dedicado a los resultados computacionales, se implementan y estudian numéricamente distintas configuraciones de PINNs aplicadas a problemas de Poisson 1D y 2D. Se presentan curvas de error en normas $L^2$ y $H^1$ en función del número de puntos de entrenamiento, comparando una PINN estándar, una penalización cuadrática media y la regularización discreta LIPR. Se examinan además los dos problemas unidimensionales de Shin et al., junto con un diagnóstico numérico de regularidad H\"older y cobertura. Por último, se comparan los resultados para Poisson 2D con métodos de elementos finitos y volúmenes finitos, dejando explícitas las diferencias entre presupuestos de cómputo y procedimientos de discretización.


